\documentclass[12pt,a4paper]{article}
\usepackage[utf-8]{inputenc}
\usepackage[T1]{fontenc}
\usepackage[portuguese]{babel}
\usepackage{geometry}
\usepackage{graphicx}
\usepackage{hyperref}
\usepackage{amsmath}
\usepackage{amssymb}
\usepackage{listings}
\usepackage{xcolor}
\usepackage{fancyhdr}
\usepackage{setspace}
\usepackage{booktabs}
\usepackage{float}

\geometry{margin=2.5cm}
\setstretch{1.5}

% Configuração de cores para código
\definecolor{codegreen}{rgb}{0,0.6,0}
\definecolor{codegray}{rgb}{0.5,0.5,0.5}
\definecolor{codepurple}{rgb}{0.58,0,0.82}
\definecolor{backcolour}{rgb}{0.95,0.95,0.92}

\lstdefinestyle{mystyle}{
    backgroundcolor=\color{backcolour},
    commentstyle=\color{codegreen},
    keywordstyle=\color{codepurple},
    numberstyle=\tiny\color{codegray},
    stringstyle=\color{codepurple},
    basicstyle=\ttfamily\footnotesize,
    breakatwhitespace=false,
    breaklines=true,
    captionpos=b,
    keepspaces=true,
    numbers=left,
    numbersep=5pt,
    showspaces=false,
    showstringspaces=false,
    showtabs=false,
    tabsize=2
}

\lstset{style=mystyle}

\title{\textbf{Sistema de Recomendação de Perfumes Nacionais} \\ 
       \large Projeto 1 --- Introdução à Inteligência Artificial}
\author{}
\date{Universidade de Brasília \\ Departamento de Ciência da Computação \\ Turma 01, 2026/1 \\ Prof. Díbio}

\pagestyle{fancy}
\fancyhf{}
\cfoot{\thepage}

\begin{document}

\maketitle

\begin{abstract}
Este projeto implementa um sistema híbrido de recomendação para uma loja virtual de perfumes nacionais, combinando as técnicas de \textbf{TF-IDF} (filtragem por conteúdo) e \textbf{SVD} (filtragem colaborativa). O sistema recomenda produtos personalizados para novos usuários a partir de suas preferências declaradas, refinadas pela matriz de utilidade de 500 usuários similares. A abordagem híbrida resolve o problema do ``cold-start'', oferecendo recomendações relevantes mesmo para usuários sem histórico de compras prévias.
\end{abstract}

\section{Introdução}

O objetivo deste projeto é desenvolver um sistema inteligente de recomendação que combine técnicas de filtragem por conteúdo e filtragem colaborativa. A motivação é prática: melhorar a experiência do usuário em uma loja virtual de perfumes nacionais, oferecendo recomendações personalizadas e relevantes.

A abordagem híbrida foi escolhida para mitigar os problemas conhecidos:
\begin{itemize}
    \item \textbf{Cold-start com TF-IDF puro:} novos usuários recebem recomendações genéricas
    \item \textbf{Cold-start com SVD puro:} sem dados de interação, não há como fazer predições
    \item \textbf{Solução híbrida:} TF-IDF gera candidatos relevantes, SVD reordena com insights colaborativos
\end{itemize}

\section{Definição do Problema e Dados}

\subsection{Catálogo de Produtos}

O projeto utiliza \textbf{50 perfumes nacionais} provenientes de 6 marcas principais:

\begin{table}[H]
\centering
\begin{tabular}{|c|c|}
\toprule
\textbf{Marca} & \textbf{Quantidade} \\
\midrule
O Boticário & 15 \\
Natura & 15 \\
Eudora & 4 \\
Avon & 7 \\
Mahogany, Granado, Phebo & 9 \\
\bottomrule
\end{tabular}
\caption{Distribuição de perfumes por marca}
\label{tab:marcas}
\end{table}

\subsubsection{Características dos Produtos}

Cada perfume possui \textbf{3 características principais} utilizadas pelo modelo de recomendação:

\begin{enumerate}
    \item \textbf{Família Olfativa} (23 tipos): Exemplos include floral, amadeirado, citríco, oriental, gourmand, aquático, aromatico, entre outros.
    
    \item \textbf{Ocasião de Uso}: Contexto de uso categorizando-se em diurno/noturno, formal/casual, verão/inverno. Exemplos: ``diurno-casual'', ``noturno-formal'', ``diurno-verao''.
    
    \item \textbf{Gênero}: Categoria de público (masculino, feminino, ou implicitamente unissex).
\end{enumerate}

Adicionalmente, cada produto inclui: preço (R\$ 69,90 a R\$ 319,90), notas olfativas (pirâmide olfativa em 5-7 componentes) e marca.

\subsection{Matriz de Utilidade}

A matriz de utilidade foi construída com \textbf{500 usuários} e \textbf{50 perfumes}, totalizando 25.000 células potenciais.

\subsubsection{Características}

\begin{itemize}
    \item \textbf{Avaliações registradas:} 5.905 (23,6\% de densidade)
    \item \textbf{Média por usuário:} 11,8 avaliações
    \item \textbf{Escala de notas:} 1 a 5 (inteiros)
    \item \textbf{Células vazias (0):} 19.095 (usuários que não avaliaram)
\end{itemize}

\subsubsection{Metodologia de Geração}

A matriz foi gerada sinteticamente mantendo critérios realistas:

\paragraph{Personas de Usuários:}
Cada usuário recebe uma ``persona'' com preferências aleatórias:
\begin{itemize}
    \item 1-2 famílias olfativas favoritas
    \item Uma ocasião preferida
    \item Conjunto de gêneros aceitáveis (45\% masc, 45\% fem, 10\% ambos)
\end{itemize}

\paragraph{Modelo de Notas:}
A nota de um perfume para um usuário é calculada como:
\begin{equation}
\text{nota} = 2.5 + \Delta_{\text{fam}} + \Delta_{\text{ocas}} + \Delta_{\text{gênero}} + \epsilon
\end{equation}

onde:
\begin{itemize}
    \item $\Delta_{\text{fam}} \in [0, 1.6]$ se há sobreposição de tokens na família olfativa
    \item $\Delta_{\text{ocas}} \in [0, 0.6]$ se há sobreposição na ocasião
    \item $\Delta_{\text{gênero}} = -2.2$ se gênero é incompatível
    \item $\epsilon \sim \mathcal{N}(0, 0.45)$ é ruído gaussiano
    \item nota final: $\text{clamp}(\text{nota}, 1, 5)$
\end{itemize}

\paragraph{Probabilidade de Avaliação:}
\begin{itemize}
    \item Gênero incompatível: 4\%
    \item Família favorita: 65\%
    \item Outros produtos: 22\%
\end{itemize}

\paragraph{Usuários Silenciosos:}
Aproximadamente 12\% dos usuários nunca deixam avaliações, modelando comportamento realista.

\section{Metodologia de Recomendação}

\subsection{Fase 1: TF-IDF (Filtragem por Conteúdo)}

O algoritmo TF-IDF identifica candidatos iniciais baseado no perfil do novo usuário.

\subsubsection{Processo}

\begin{enumerate}
    \item Construir vetor de características do usuário: $\text{perfil} = \text{família} + \text{ ocasião} + \text{ gênero}$
    
    \item Vectorizar usando TF-IDF pré-treinado:
    \begin{equation}
    \vec{v}_{\text{perfil}} = \text{TfidfVectorizer.transform}(\text{perfil})
    \end{equation}
    
    \item Calcular similaridade de cosseno com cada perfume:
    \begin{equation}
    \text{sim}(i) = \frac{\vec{v}_{\text{perfil}} \cdot \vec{v}_{i}}{|\vec{v}_{\text{perfil}}| \cdot |\vec{v}_{i}|}
    \end{equation}
    
    \item Aplicar filtros:
    \begin{itemize}
        \item Gênero deve corresponder (ou expandir se nenhum resultado)
        \item Preço $\leq$ limite do usuário
    \end{itemize}
    
    \item Retornar top-20 candidatos ordenados por $\text{sim}(i)$
\end{enumerate}

\subsubsection{Características}

\begin{itemize}
    \item \textbf{Velocidade:} O(n) em número de perfumes
    \item \textbf{Funciona com novos usuários:} Não requer histórico
    \item \textbf{Determinístico:} Mesmo perfil $\Rightarrow$ mesmas recomendações
    \item \textbf{Transparente:} Fácil explicar ``por que'' um perfume foi recomendado
\end{itemize}

\subsection{Fase 2: SVD (Filtragem Colaborativa)}

O algoritmo SVD reordena os top-20 candidatos usando matrix factorization e insights de usuários similares.

\subsubsection{Processo}

\begin{enumerate}
    \item Encontrar os 50 usuários mais similares ao novo usuário usando scores TF-IDF
    
    \item Extrair predições do SVD pré-treinado para esses 50 usuários:
    \begin{equation}
    \hat{R}_{50 \times 20} = \text{pred\_matrix}[\text{top\_50\_usuarios}, \text{top\_20\_perfumes}]
    \end{equation}
    
    \item Calcular nota média prevista por perfume:
    \begin{equation}
    \text{svd\_score}(j) = \frac{1}{50} \sum_{i=1}^{50} \hat{R}_{ij}
    \end{equation}
    
    \item Normalizar para [0, 1]:
    \begin{equation}
    \text{svd\_norm}(j) = \frac{\text{svd\_score}(j) - 1}{4}
    \end{equation}
    
    \item Combinar com TF-IDF (vide seção 3.3)
\end{enumerate}

\subsubsection{Características}

\begin{itemize}
    \item Captura correlações ocultas entre produtos
    \item Beneficia-se de padrões coletivos
    \item Descoberta de serendipidade (usuário pode gostar de perfume não óbvio)
\end{itemize}

\subsection{Combinação Híbrida}

O score final combina ambas as abordagens:

\begin{equation}
\text{score\_final}(j) = 0.7 \times \text{tfidf}(j) + 0.3 \times \text{svd\_norm}(j)
\end{equation}

\textbf{Justificativa dos pesos:}
\begin{itemize}
    \item 70\% TF-IDF: Mantém relevância de conteúdo como base principal
    \item 30\% SVD: Incorpora insights colaborativos de forma complementar
    \item Evita overfitting a usuários similares (podem ter preferências atípicas)
\end{itemize}

O sistema retorna os \textbf{top-5} perfumes ordenados por score\_final.

\section{Implementação}

\subsection{Estrutura do Projeto}

\begin{lstlisting}[language=bash]
sistema-recomendacao/
├── app.py                      # Interface Gradio (UI web)
├── scripts/
│   ├── gerar_dados.py          # Gera catálogo + matriz (seed=42)
│   └── sistema_recomendacao.py # Código modularizado do notebook
├── notebook/
│   └── sistema_recomendacao.ipynb
├── modelos/                    # Pré-treinados
│   ├── tfidf_vectorizer.pkl
│   ├── tfidf_matrix.pkl
│   └── svd_pred_matrix.pkl
└── dados/
    ├── produtos.csv
    ├── produtos_processados.csv
    ├── matriz_utilidade.csv
    └── feedback.csv
\end{lstlisting}

\subsection{Dependências}

\begin{itemize}
    \item \textbf{pandas}: manipulação de dados
    \item \textbf{scikit-learn}: TF-IDF, SVD, similaridade de cosseno
    \item \textbf{scipy}: operações numéricas
    \item \textbf{seaborn, matplotlib}: visualizações (notebook)
    \item \textbf{gradio}: interface web
    \item \textbf{joblib}: serialização de modelos
\end{itemize}

\subsection{Interface Gradio}

A interface web oferece 4 abas funcionais:

\paragraph{1. Portal do Cliente}
\begin{itemize}
    \item Cadastro: email, senha, nome, preferências
    \item Login/Logout
    \item Editar preferências (accordion)
\end{itemize}

\paragraph{2. Minha Vitrine}
\begin{itemize}
    \item Botão ``Gerar Recomendações''
    \item Exibe top-5 com imagem, preço, família olfativa
    \item Adiciona ao carrinho e finaliza pedidos
\end{itemize}

\paragraph{3. Feedback}
\begin{itemize}
    \item Carrega pedidos realizados
    \item Usuário avalia cada perfume (1-5)
    \item Salva em \texttt{dados/feedback.csv} (persistido)
\end{itemize}

\paragraph{4. Portal Anônimo}
Demo sem necessidade de cadastro para explorar o sistema.

\section{Resultados}

\subsection{Matriz de Utilidade}

\begin{itemize}
    \item 500 usuários, 50 perfumes, 25.000 células
    \item 5.905 avaliações (23,6\% de densidade)
    \item 11,8 avaliações por usuário (em média)
    \item Densidade adequada para SVD (recomenda-se > 10\%)
\end{itemize}

\subsection{Desempenho Computacional}

\begin{table}[H]
\centering
\begin{tabular}{|l|r|}
\toprule
\textbf{Tarefa} & \textbf{Tempo} \\
\midrule
Geração de dados & < 1s \\
Treinamento TF-IDF & < 1s \\
Treinamento SVD & < 2s \\
Recomendação por usuário & < 100ms \\
\bottomrule
\end{tabular}
\caption{Tempos de execução}
\label{tab:performance}
\end{table}

\subsection{Qualidade das Recomendações}

O sistema prioriza:
\begin{enumerate}
    \item Relevância de conteúdo (família olfativa, ocasião, gênero)
    \item Preferências colaborativas (usuários similares)
    \item Filtros práticos (gênero, faixa de preço)
\end{enumerate}

\textbf{Exemplo de fluxo:}
\begin{itemize}
    \item \textit{Entrada:} Nova usuária, feminino, gosta de florais, orçamento R\$ 150
    \item \textit{TF-IDF (top-20):} Lily Essence, Floratta Rosas, Egeo Dolce, \ldots
    \item \textit{SVD reordena:} Baseado em padrões de mulheres com preferência floral
    \item \textit{Top-5 final:} Perfumes florais bem-avaliados por personas similares
\end{itemize}

\subsection{Limitações e Tratamento}

\paragraph{Cold-start para novos usuários:}
Sem histórico, SVD retorna predições neutras. Nesse caso, TF-IDF (70\%) domina, garantindo recomendações minimamente relevantes. À medida que o usuário interage, seria possível retreinar SVD periodicamente.

\paragraph{Matriz esparsa:}
23,6\% de densidade é adequado para SVD, mas matriz mais densa melhoraria predições.

\section{Conclusão}

O projeto implementa com sucesso um sistema de recomendação híbrido que atende todos os requisitos:

\begin{itemize}
    \item[$\checkmark$] 50 produtos com 3 características + avaliação 1-5
    \item[$\checkmark$] 500 usuários com matriz de utilidade reproduzível (seed=42)
    \item[$\checkmark$] Modelo híbrido TF-IDF + SVD bem documentado
    \item[$\checkmark$] Interface web funcional (Gradio)
    \item[$\checkmark$] Notebook Jupyter com EDA e treinamento
    \item[$\checkmark$] Código modularizado e reutilizável
\end{itemize}

A abordagem híbrida oferece um balanço entre precisão de conteúdo (TF-IDF) e descoberta de padrões colaborativos (SVD), resultando em recomendações relevantes e explicáveis mesmo para usuários sem histórico.

\subsection{Extensões Futuras}

\begin{enumerate}
    \item Incorporar feedback do usuário ao retreinar SVD periodicamente
    \item Usar deep learning (autoencoders, embeddings) para representações contínuas
    \item A/B testing de pesos (variar 0.7 e 0.3)
    \item Análise de serendipidade e cobertura
    \item Recomendações personalizadas por contexto (sazonalidade, clima)
\end{enumerate}

\section*{Referências}

\begin{itemize}
    \item Ricci, F., Rokach, L., \& Shapira, B. (Eds.). (2015). \textit{Recommender Systems Handbook}. Springer.
    \item Lops, P., De Gemmis, M., \& Semeraro, G. (2011). Content-based recommender systems: State of the art and trends. In \textit{Recommender systems handbook} (pp. 73-105). Springer.
    \item Koren, Y., Bell, R., \& Volinsky, C. (2009). Matrix factorization techniques for recommender systems. \textit{IEEE computer}, 42(8), 30-37.
\end{itemize}

\end{document}
